\documentclass[a4paper,11pt]{article}
\usepackage[utf8]{inputenc}
\usepackage[czech]{babel}
\usepackage{amsmath, amssymb, amsthm}
\usepackage[margin=2cm]{geometry}
\usepackage{enumerate}

\title{\textbf{Velká písemná práce z Lineární algebry 1}}
\author{Zimní semestr -- Soustavy až Determinanty}
\date{}

\begin{document}

\maketitle

\noindent \textbf{Instrukce:}
\begin{itemize}
    \item Celkový počet bodů: 100.
    \item Pro zisk zápočtu je třeba alespoň 60 bodů.
    \item \textbf{Sekce I (Ano/Ne):} K získání bodů za otázku je třeba odpovědět správně na všechny tři její podčásti.
\end{itemize}

\hrule
\vspace{0.5cm}

\section*{Část I: Ano/Ne (8 bodů)}
\textit{Rozhodněte, zda platí následující tvrzení. Není třeba zdůvodňovat. (2 body za každou správně zodpovězenou trojici)}

\begin{enumerate}[(a)]
    \item \textbf{Soustavy rovnic:} Uvažme soustavu lineárních rovnic $Ax=b$ nad $\mathbb{R}$ o $m$ rovnicích a $n$ neznámých, kde $m < n$.
    \begin{itemize}
        \item[\textbf{ANO / NE}] Soustava má vždy alespoň jedno řešení.
        \item[\textbf{ANO / NE}] Pokud má soustava alespoň jedno řešení, pak jich má nekonečně mnoho.
        \item[\textbf{ANO / NE}] Množina všech řešení této soustavy tvoří podprostor vektorového prostoru $\mathbb{R}^n$ právě tehdy, když $b$ je nulový vektor.
    \end{itemize}

    \item \textbf{Lineární závislost:} Nechť $u, v, w$ jsou lineárně nezávislé vektory ve vektorovém prostoru $V$ nad tělesem $T$.
    \begin{itemize}
        \item[\textbf{ANO / NE}] Vektory $u, v$ jsou lineárně nezávislé.
        \item[\textbf{ANO / NE}] Vektory $u+v, v+w, w+u$ jsou lineárně nezávislé (předpokládejme $\text{char}(T) \neq 2$).
        \item[\textbf{ANO / NE}] Pokud $z \in \text{span}\{u, v, w\}$, pak je množina $\{u, v, w, z\}$ lineárně nezávislá.
    \end{itemize}

    \item \textbf{Matice a zobrazení:} Nechť $A$ je čtvercová matice řádu $n$ nad $\mathbb{R}$.
    \begin{itemize}
        \item[\textbf{ANO / NE}] Pokud $A^2 = I$ (jednotková matice), pak je $A$ regulární.
        \item[\textbf{ANO / NE}] Pokud je rovnice $Ax=0$ řešitelná pouze triviálně (tj. $x=0$), pak je zobrazení $f(x) = Ax$ na (surjektivní).
        \item[\textbf{ANO / NE}] Hodnost matice $A$ se může změnit, pokud k jednomu řádku přičteme násobek jiného řádku.
    \end{itemize}

    \item \textbf{Determinanty:} Nechť $A, B$ jsou čtvercové matice řádu $n$ nad $\mathbb{R}$.
    \begin{itemize}
        \item[\textbf{ANO / NE}] Pro libovolné matice platí $\det(A+B) = \det(A) + \det(B)$.
        \item[\textbf{ANO / NE}] Platí $\det(AB) = \det(BA)$.
        \item[\textbf{ANO / NE}] Pokud je matice $A$ singulární, pak $\det(A^T) = 0$.
    \end{itemize}
\end{enumerate}

\vspace{0.5cm}

\section*{Část II: Definice (12 bodů)}
\textit{Uveďte přesnou definici následujících pojmů (včetně předpokladů). (3 body za každou definici)}

\begin{enumerate}
    \item Lineární obal množiny vektorů $M$ ve vektorovém prostoru $V$.
    \item Báze vektorového prostoru.
    \item Matice přechodu od báze $B$ k bázi $C$.
    \item Determinant čtvercové matice (pomocí permutací nebo rozvoje).
\end{enumerate}

\vspace{0.5cm}

\section*{Část III: Jednoduché příklady (15 bodů)}
\textit{Vypočtěte nebo určete. Stačí uvést správný výsledek. (3 body za každý příklad)}

\begin{enumerate}
    \item Určete hodnost matice $A = \begin{pmatrix} 1 & 2 & 1 \\ 2 & 4 & 2 \\ -1 & -2 & 0 \end{pmatrix}$ nad tělesem $\mathbb{R}$.
    \item Najděte souřadnice vektoru $v = (3, 1)^T$ v bázi $B = \{ (1, 1)^T, (0, 1)^T \}$.
    \item Vypočtěte součin permutací $\pi \circ \sigma$ v $S_5$, kde $\pi = (1, 3, 4)$ a $\sigma = (1, 2, 5)$. Zapište výsledek pomocí nezávislých cyklů.
    \item Pro jaké $a \in \mathbb{R}$ je matice $\begin{pmatrix} 1 & a \\ a & 1 \end{pmatrix}$ singulární?
    \item Jak se změní determinant matice $A$ typu $3 \times 3$, pokud vynásobíme její první řádek číslem 3 a druhý řádek číslem 2?
\end{enumerate}

\vspace{0.5cm}

\section*{Část IV: Formulace tvrzení (12 bodů)}
\textit{Zformulujte přesně následující věty/tvrzení (není třeba dokazovat). (3 body za každé tvrzení)}

\begin{enumerate}
    \item Frobeniova věta (o řešitelnosti soustav lineárních rovnic).
    \item Steinitzova věta o výměně.
    \item Věta o dimenzi jádra a obrazu lineárního zobrazení.
    \item Cauchyova-Binetova věta (o determinantu součinu matic).
\end{enumerate}

\newpage

\section*{Část V: Početní příklady s postupem (15 bodů)}
\textit{Vyřešte následující úlohy. Uveďte podrobný postup řešení. (5 bodů za každý příklad)}

\begin{enumerate}
    \item Najděte množinu všech řešení soustavy lineárních rovnic nad $\mathbb{R}$ zadanou rozšířenou maticí:
    $$
    \left( \begin{array}{ccc|c}
    1 & 2 & 3 & 4 \\
    2 & 5 & 8 & 9 \\
    3 & 6 & 10 & 10
    \end{array} \right)
    $$
    
    \item K dané matici $A$ nad tělesem $\mathbb{Z}_5$ nalezněte matici inverzní, nebo ukažte, že neexistuje:
    $$
    A = \begin{pmatrix} 1 & 2 & 4 \\ 2 & 0 & 1 \\ 3 & 1 & 2 \end{pmatrix}
    $$
    
    \item Vypočítejte determinant matice $4 \times 4$ nad $\mathbb{R}$ převodem na odstupňovaný tvar nebo rozvojem:
    $$
    D = \det \begin{pmatrix} 
    2 & 1 & 0 & 0 \\
    1 & 2 & 1 & 0 \\
    0 & 1 & 2 & 1 \\
    0 & 0 & 1 & 2 
    \end{pmatrix}
    $$
\end{enumerate}

\vspace{0.5cm}

\section*{Část VI: Důkazy jednodušších tvrzení (20 bodů)}
\textit{Dokažte následující tvrzení. (5 bodů za každý důkaz)}

\begin{enumerate}
    \item Nechť $A, B$ jsou čtvercové regulární matice stejného řádu. Dokažte, že $(AB)^{-1} = B^{-1}A^{-1}$.
    \item Dokažte, že průnik dvou podprostorů $U, V$ vektorového prostoru $W$ je opět podprostor $W$.
    \item Nechť $f: U \to V$ je lineární zobrazení. Dokažte: Zobrazení $f$ je prosté (injektivní) právě tehdy, když $\ker(f) = \{o\}$.
    \item Dokažte, že pro libovolnou matici $A$ platí, že elementární řádkové úpravy nemění lineární obal řádků matice $A$.
\end{enumerate}

\vspace{0.5cm}

\section*{Část VII: Úlohy na zamyšlení (18 bodů)}
\textit{Vyřešte nebo dokažte následující náročnější úlohy. (6 bodů za každou úlohu)}

\begin{enumerate}
    \item \textbf{Důkaz hodnosti:} Nechť $A$ je matice typu $m \times n$ a $B$ je matice typu $n \times p$. Dokažte, že hodnost součinu $AB$ není větší než hodnost matice $A$, tj. $rank(AB) \leq rank(A)$.
    
    \item \textbf{Prostor polynomů:} Uvažujme vektorový prostor $P_2$ všech polynomů stupně nejvýše 2 nad $\mathbb{R}$. Rozhodněte (a dokažte), zda jsou polynomy $p_1(x) = x^2 - x + 1$, $p_2(x) = x^2 + 2x$ a $p_3(x) = 3x - 1$ lineárně nezávislé. Tvoří bázi $P_2$?
    
    \item \textbf{Vandermondův determinant:}
    Nechť $x_1, x_2, x_3 \in \mathbb{R}$. Dokažte výpočtem determinantu, že matice
    $$
    V = \begin{pmatrix}
    1 & x_1 & x_1^2 \\
    1 & x_2 & x_2^2 \\
    1 & x_3 & x_3^2
    \end{pmatrix}
    $$
    je regulární právě tehdy, když jsou čísla $x_1, x_2, x_3$ navzájem různá.
\end{enumerate}

\end{document}